\documentclass[
    12pt,                % tamanho da fonte
    openright,           % capítulos começam em página ímpar
    oneside,             % para impressão em apenas um lado do papel
    a4paper,             % tamanho do papel
    english,             % idioma adicional
    brazil               % idioma principal
]{abntex2}

% ---------------------------------------------------------------
% PACOTES
% ---------------------------------------------------------------
\usepackage{lmodern}                % Fonte Latin Modern
\usepackage[T1]{fontenc}            % Seleção de códigos de fonte
\usepackage[utf8]{inputenc}         % Codificação do documento
\usepackage{lastpage}               % Usado pela Ficha catalográfica
\usepackage{indentfirst}            % Indenta o primeiro parágrafo
\usepackage{color}                  % Controle de cores
\usepackage{graphicx}               % Inclusão de gráficos
\usepackage{microtype}              % Melhorias de justificação
\usepackage{float}                  % Para posicionamento de figuras
\usepackage{booktabs}               % Tabelas profissionais
\usepackage{multirow}               % Células mescladas em tabelas
\usepackage{longtable}              % Tabelas longas
\usepackage{amsmath,amssymb}        % Símbolos matemáticos
\usepackage{pgfgantt}               % Diagramas de Gantt
\usepackage[brazilian,hyperpageref]{backref}  % Páginas com citações
\usepackage[alf]{abntex2cite}       % Citações ABNT

% ---------------------------------------------------------------
% CONFIGURAÇÕES DO PDF
% ---------------------------------------------------------------
\makeatletter
\hypersetup{
    pdftitle={\@Classificação de cancer de mama utilizando Redes Neurais Convolucionais clássicas e quânticas},
    pdfauthor={\@Matheus-Sanchez-Duda},
    pdfsubject={TCC1 - SENAC Santo Amaro},
    pdfcreator={LaTeX with abnTeX2},
    pdfkeywords={tcc}{senac}{santo amaro}{trabalho de conclusão},
    colorlinks=true,        % false: links em caixas; true: links coloridos
    linkcolor=blue,         % cor dos links internos
    citecolor=blue,         % cor dos links para bibliografia
    filecolor=magenta,      % cor dos links para arquivos
    urlcolor=blue,
    bookmarksdepth=4
}
\makeatother

% ---------------------------------------------------------------
% CONFIGURAÇÕES DE APARÊNCIA DO PDF FINAL
% ---------------------------------------------------------------
\definecolor{blue}{RGB}{41,5,195}

% ---------------------------------------------------------------
% INFORMAÇÕES DO DOCUMENTO
% ---------------------------------------------------------------
\titulo{Título do Trabalho de Conclusão de Curso}
\autor{Nome Completo do Aluno}
\local{São Paulo}
\data{2025}
\orientador{Prof. Nome do Orientador}
% \coorientador{Prof. Nome do Coorientador} % Descomente se houver coorientador


\instituicao{%
  SENAC -- Serviço Nacional de Aprendizagem Comercial
  \par
  Unidade Santo Amaro
  \par
  Curso de [Nome do Curso]
}

\tipotrabalho{Trabalho de Conclusão de Curso I}

\preambulo{Proposta de Trabalho de Conclusão de Curso apresentado ao SENAC Santo Amaro como requisito parcial para aprovação na disciplina de TCC1 do curso de [Nome do Curso].}

% ---------------------------------------------------------------
% COMPILA O ÍNDICE
% ---------------------------------------------------------------
\makeindex

% ---------------------------------------------------------------
% INÍCIO DO DOCUMENTO
% ---------------------------------------------------------------
\begin{document}

% Seleciona o idioma do documento
\selectlanguage{brazil}

% Retira espaço extra obsoleto entre as frases
\frenchspacing

% ---------------------------------------------------------------
% ELEMENTOS PRÉ-TEXTUAIS
% ---------------------------------------------------------------
\pretextual

% ---------------------------------------------------------------
% CAPA
% ---------------------------------------------------------------
\imprimircapa

% ---------------------------------------------------------------
% FOLHA DE ROSTO
% ---------------------------------------------------------------
\imprimirfolhaderosto

% ---------------------------------------------------------------
% RESUMO
% ---------------------------------------------------------------
\begin{resumo}
Este documento apresenta a proposta inicial do Trabalho de Conclusão de Curso (TCC1) desenvolvido no SENAC Santo Amaro. O objetivo deste trabalho é [descrever brevemente o objetivo principal do TCC]. A justificativa para este estudo baseia-se em [descrever brevemente a justificativa]. A revisão bibliográfica abrange [mencionar os principais tópicos]. Como solução proposta, pretende-se [descrever brevemente a solução]. Este trabalho está organizado em três capítulos que apresentam a introdução ao tema, a revisão bibliográfica que fundamenta a pesquisa, o desenvolvimento com a solução proposta e o cronograma de trabalho, seguidos das referências bibliográficas.

\vspace{\onelineskip}

\noindent
\textbf{Palavras-chave}: palavra1. palavra2. palavra3. palavra4. palavra5.
\end{resumo}

% ---------------------------------------------------------------
% LISTA DE ILUSTRAÇÕES
% ---------------------------------------------------------------
% \pdfbookmark[0]{\listfigurename}{lof}
% \listoffigures*
% \cleardoublepage

% ---------------------------------------------------------------
% LISTA DE TABELAS
% ---------------------------------------------------------------
% \pdfbookmark[0]{\listtablename}{lot}
% \listoftables*
% \cleardoublepage

% ---------------------------------------------------------------
% SUMÁRIO
% ---------------------------------------------------------------
\pdfbookmark[0]{\contentsname}{toc}
\tableofcontents*
\cleardoublepage

% ---------------------------------------------------------------
% ELEMENTOS TEXTUAIS
% ---------------------------------------------------------------
\textual

